\subsection*{Câu 1}

\noindent\textbf{1. Mạch kiến thức:} Thứ tự thực hiện phép tính

\noindent\textbf{2. Độ khó:} *

\noindent\textbf{3. Đề bài:}
\[
10 + 36 : 2 \cdot 3
\]

\noindent\textbf{4. Lời giải:}
\begin{align*}
10 + 36 : 2 \cdot 3
    &= 46 : 2 \cdot 3 \\
    &= 23 \cdot 3 \\
    &= 69
\end{align*}

\noindent\textbf{5. Token bắt đầu sai:}
\[
46
\]

\noindent\textbf{6. Giải thích lỗi:}

Theo thứ tự thực hiện phép tính, phép nhân và phép chia được thực hiện trước phép cộng và phép trừ.

\noindent\textbf{7. Sửa lại:}

\textit{Dạng câu hỏi:} Tự luận

\textit{Câu hỏi:} Viết lại phép biến đổi cho đúng.

\textit{Đáp án:}
\[
10 + 18 * 3
\]

\subsection*{Câu 2}

\noindent\textbf{1. Mạch kiến thức:} Quy tắc dấu ngoặc

\noindent\textbf{2. Độ khó:} *

\noindent\textbf{3. Đề bài:}
\[
214-(-36)+(-305)
\]

\noindent\textbf{4. Lời giải:}
\begin{align*}
214-(-36)+(-305)
    &= 214-36-305 \\
    &= 178-305 \\
    &= -127
\end{align*}

\noindent\textbf{5. Token bắt đầu sai:}
\[
-
\]

\noindent\textbf{6. Giải thích lỗi:}\\
Khi bỏ dấu ngoặc có dấu ``$-$'' đằng trước, phải đổi dấu của tất cả các số hạng trong dấu ngoặc.

\noindent\textbf{7. Sửa lại:}

\textit{Dạng câu hỏi:} Trắc nghiệm

\textit{Câu hỏi:} Dấu ``$-$'' trong phép biến đổi trên cần được sửa thành dấu nào?

\[
214 \boxed{\phantom{+}}36 - 305
\]

\begin{itemize}
    \item[A.] $+$
    \item[B.] $-$
    \item[C.] $*$
    \item[D.] $:$
\end{itemize}

\textit{Đáp án:}
\[
A. +
\]

\subsection*{Câu 3}

\noindent\textbf{1. Mạch kiến thức:} Kết hợp thứ tự thực hiện phép tính và quy tắc dấu ngoặc

\noindent\textbf{2. Độ khó:} **

\noindent\textbf{3. Đề bài:}
\[
1+2*(2^3-2*3)
\]

\noindent\textbf{4. Lời giải:}
\begin{align*}
1+2*(2^3-2*3)
    &= 1+2*(8-2*3) \\
    &= 1+2*(6*3) \\
    &= 1+2*18 \\
    &= 1+36 \\
    &= 37
\end{align*}

\noindent\textbf{5. Token bắt đầu sai:}
\[
6
\]

\noindent\textbf{6. Giải thích lỗi:}

Thực hiện trong dấu ngoặc trước; trong ngoặc: lũy thừa $\rightarrow$ nhân, chia $\rightarrow$ cộng, trừ.

\noindent\textbf{7. Sửa lại:}

\textit{Dạng câu hỏi:} Tự luận

\textit{Câu hỏi:} Viết lại phép biến đổi cho đúng.

\textit{Đáp án:}
\[
1+2*(8-6)
\]

\subsection*{Câu 4}

\noindent\textbf{1. Mạch kiến thức:} Thứ tự thực hiện phép tính

\noindent\textbf{2. Độ khó:} **

\noindent\textbf{3. Đề bài:}
\[
1,2 - 3^2 + 7,5 : 3
\]

\noindent\textbf{4. Lời giải:}
\begin{align*}
1,2 - 3^2 + 7,5 : 3
    &= 1,2 - 9 + 7,5 : 3 \\
    &= -7,7 + 7,5 : 3 \\
    &= -7,7 + 2,5 \\
    &= -5,2
\end{align*}

\noindent\textbf{5. Token bắt đầu sai:}
\[
-7,7
\]

\noindent\textbf{6. Giải thích lỗi:}

Nếu biểu thức có các phép tính cộng, trừ, nhân, chia và lũy thừa thì thực hiện theo thứ tự: Lũy thừa $\rightarrow$ Nhân, chia $\rightarrow$ Cộng, trừ.

\noindent\textbf{7. Sửa lại:}

\textit{Dạng câu hỏi:} Tự luận

\textit{Câu hỏi:} Viết lại phép biến đổi cho đúng:$1,2 - 9 + 7,5 : 3=...$ 

\textit{Đáp án:}
\[
1,2 - 9 + 2.5
\]

\subsection*{Câu 5}

\noindent\textbf{1. Mạch kiến thức:} Quy tắc chuyển vế

\noindent\textbf{2. Độ khó:} **

\noindent\textbf{3. Đề bài:}
\[
-\frac{1}{3}-x=\frac{17}{6}
\]

\noindent\textbf{4. Lời giải:}
\begin{align*}
-\frac{1}{3}-x
    &= \frac{17}{6} \\
x
    &= -\frac{1}{3}+\frac{17}{6} \\
x
    &= -\frac{2}{6}+\frac{17}{6}
    \\
x
    &= \frac{5}{2}
\end{align*}

\noindent\textbf{5. Token bắt đầu sai:}
\[
+
\]

\noindent\textbf{6. Giải thích lỗi:}

Trong một đẳng thức, khi chuyển một số hạng từ vế này sang vế kia, ta phải đổi dấu của số hạng đó.

\noindent\textbf{7. Sửa lại:}

\textit{Dạng câu hỏi:} Trắc nghiệm

\textit{Câu hỏi:} Dấu ``$+$'' trong phép biến đổi trên cần được sửa thành dấu nào?

\[
x=-\frac{1}{3}\ \boxed{\phantom{+}}\ \frac{17}{6}
\]

\begin{itemize}
    \item[A.] $+$
    \item[B.] $-$
    \item[C.] $*$
    \item[D.] $:$
\end{itemize}

\textit{Đáp án:} B